\documentclass[11pt]{article}
\usepackage[margin=1in]{geometry}
\usepackage{booktabs}
\usepackage{array}
\usepackage{longtable}
\usepackage{hyperref}
\usepackage[T1]{fontenc}
\usepackage[utf8]{inputenc}
\usepackage{lmodern}

\title{Qwen3 Loop-Trigger Attention Query Comparison}
\author{}
\date{2026-04-14}

\begin{document}
\maketitle

\section*{Object}

This note compares the same full Qwen3 loop-trigger rerun at two query positions on the saved March 22--23 prompt-profile archives:

\begin{itemize}
  \item \texttt{trigger\_end}: the final token of the twentieth repeated $30$-gram
  \item \texttt{trigger\_start}: the first token of that same twentieth repeated $30$-gram
\end{itemize}

Everything else is held fixed. The analyzed object is still:

\begin{itemize}
  \item model: \texttt{Qwen/Qwen3-1.7B}
  \item loop definition: \texttt{n=30}, \texttt{k=20}
  \item selected rows: \texttt{811 / 820} replayable loop rows
  \item dataset counts: \texttt{GPQA 121}, \texttt{AIME 30}, \texttt{MATH-500 68}, \texttt{MMLU-Pro 130}, \texttt{LiveCodeBench 462}
  \item total prefix lengths: median \texttt{9846}, \texttt{p95 = 22321}, max \texttt{28830}
\end{itemize}

There is no extra mid-prefix truncation inside either run. Each HF forward still sees the full prompt plus the completion prefix through the saved first loop trigger. Only the query token changes.

\section*{Reconstruction Boundary}

The March prompt-profile archives preserve exact \texttt{prompt\_token\_ids}, but they do not preserve exact \texttt{completion\_token\_ids}. For those older rows, the replay path still has to retokenize saved \texttt{completion\_text}.

\begin{center}
\small
\begin{tabular}{lrrrrr}
\toprule
Dataset & Rollouts & Loop rows & Exact length & +1 hidden stop & Trigger-prefix match \\
\midrule
GPQA & 792 & 124 & 46 / 792 & 749 / 792 & 121 / 124 \\
AIME & 240 & 31 & 38 / 240 & 234 / 240 & 30 / 31 \\
MATH-500 & 2000 & 68 & 38 / 2000 & 1991 / 2000 & 68 / 68 \\
MMLU-Pro & 3200 & 130 & 34 / 3200 & 3191 / 3200 & 130 / 130 \\
LiveCodeBench & 3200 & 467 & 356 / 3200 & 3168 / 3200 & 462 / 467 \\
\midrule
Overall & 9432 & 820 & 512 / 9432 & 9333 / 9432 & 811 / 820 \\
\bottomrule
\end{tabular}
\normalsize
\end{center}

So both analyses remain exact prompt replay plus approximate completion replay. The important boundary is unchanged: the trigger prefix is recovered on \texttt{811 / 820} loop rows.

\section*{How The Measurement Is Defined}

\begin{enumerate}
  \item Recompute the loop trigger with \texttt{find\_ngram\_loop\_trigger(..., n=30, k=20)}.
  \item Keep only rows where the recomputed \texttt{first\_loop\_prefix} matches the saved \texttt{first\_loop\_prefix\_length}.
  \item Replay the model on \texttt{prompt\_token\_ids + completion\_token\_ids[:first\_loop\_prefix\_length]}.
  \item Compare two query positions on that same replay:
  \begin{itemize}
    \item \texttt{trigger\_end = prompt\_len + trigger\_end = total\_prefix\_length - 1}
    \item \texttt{trigger\_start = prompt\_len + trigger\_start}
  \end{itemize}
  \item For every layer and head, capture the softmax attention from that query token while preserving the model's masking behavior.
  \item Partition keys into disjoint bins: \texttt{prompt}, \texttt{current\_trigger}, \texttt{previous\_loop}, \texttt{last\_previous\_loop}, \texttt{recent\_nonloop}, and \texttt{other\_completion}.
  \item Summaries are:
  \begin{itemize}
    \item \texttt{mean\_*\_mass}: bin mass averaged across heads, then across selected rows
    \item \texttt{top1\_fraction\_*}: fraction of heads whose argmax lands in that bin
  \end{itemize}
\end{enumerate}

For \texttt{trigger\_start}, the bins are still defined over the full trigger span, but the model remains causal. So keys to the right of the query token get zero mass automatically. That means \texttt{current\_trigger} at \texttt{trigger\_start} is effectively dominated by the first trigger token/self.

\section*{Final-Layer Overall Comparison}

The final-layer summary is layer $27$ for every dataset.

\begin{center}
\small
\begin{tabular}{lrrrrrrrr}
\toprule
Query & Prev-loop & Prompt & Current-trigger & Recent-nonloop & Top-1 prev & Top-1 prompt & Top-1 current & Top-1 recent \\
\midrule
\texttt{trigger\_end} & 0.031 & 0.639 & 0.176 & 0.038 & 0.0003 & 0.874 & 0.114 & 0.0005 \\
\texttt{trigger\_start} & 0.054 & 0.576 & 0.170 & 0.091 & 0.019 & 0.825 & 0.120 & 0.029 \\
\bottomrule
\end{tabular}
\normalsize
\end{center}

So moving the query from the end of the triggering copy to its first token does change the late-layer read: previous-loop mass rises from \texttt{0.031} to \texttt{0.054}, previous-loop top-1 fraction rises from essentially zero to \texttt{0.019}, prompt mass falls from \texttt{0.639} to \texttt{0.576}, and recent-nonloop mass rises from \texttt{0.038} to \texttt{0.091}. But \texttt{trigger\_start} is still prompt-dominant in the final layer.

\section*{Trigger-End Final Layer}

\begin{center}
\small
\begin{tabular}{lrrrrr}
\toprule
Dataset & Rows & Prev-loop & Prompt & Current-trigger & Recent-nonloop \\
\midrule
GPQA & 121 & 0.025 & 0.671 & 0.173 & 0.039 \\
AIME & 30 & 0.038 & 0.627 & 0.176 & 0.027 \\
MATH-500 & 68 & 0.029 & 0.634 & 0.185 & 0.039 \\
MMLU-Pro & 130 & 0.027 & 0.620 & 0.185 & 0.050 \\
LiveCodeBench & 462 & 0.034 & 0.637 & 0.174 & 0.034 \\
\midrule
Overall row-weighted & 811 & 0.031 & 0.639 & 0.176 & 0.038 \\
\bottomrule
\end{tabular}
\normalsize
\end{center}

\section*{Trigger-Start Final Layer}

\begin{center}
\small
\begin{tabular}{lrrrrr}
\toprule
Dataset & Rows & Prev-loop & Prompt & Current-trigger & Recent-nonloop \\
\midrule
GPQA & 121 & 0.035 & 0.608 & 0.171 & 0.095 \\
AIME & 30 & 0.073 & 0.564 & 0.163 & 0.073 \\
MATH-500 & 68 & 0.061 & 0.584 & 0.156 & 0.092 \\
MMLU-Pro & 130 & 0.031 & 0.565 & 0.177 & 0.120 \\
LiveCodeBench & 462 & 0.063 & 0.570 & 0.170 & 0.083 \\
\midrule
Overall row-weighted & 811 & 0.054 & 0.576 & 0.170 & 0.091 \\
\bottomrule
\end{tabular}
\normalsize
\end{center}

At \texttt{trigger\_start}, overall row-weighted top-1 fractions are:

\begin{itemize}
  \item prompt: \texttt{0.825}
  \item current trigger: \texttt{0.120}
  \item previous loop: \texttt{0.019}
  \item recent nonloop: \texttt{0.029}
\end{itemize}

So the stricter first-token query does pick up more late-layer previous-loop attention, but prompt tokens remain the dominant target on every dataset.

\section*{Mid-Stack Comparison}

Earlier loop copies remain much more visible in the middle of the stack, and the peak layer stays $6$ for every dataset under both query definitions.

\begin{center}
\small
\begin{tabular}{lrrrrrr}
\toprule
Query & Peak layer & Prev-loop & Prompt & Current-trigger & Top-1 prev & Top-1 prompt \\
\midrule
\texttt{trigger\_end} & 6 & 0.193 & 0.265 & 0.223 & 0.139 & 0.596 \\
\texttt{trigger\_start} & 6 & 0.249 & 0.245 & 0.106 & 0.197 & 0.605 \\
\bottomrule
\end{tabular}
\normalsize
\end{center}

At \texttt{trigger\_start}, the mid-stack previous-loop signal is plainly stronger than at \texttt{trigger\_end}. The clearest dataset-level peak is AIME, where layer $6$ reaches \texttt{prev-loop mass = 0.299}; the weakest is GPQA, which still reaches \texttt{0.218}.

\section*{Conclusion}

The updated read is:

\begin{itemize}
  \item the strong claim ``Qwen3 is not paying attention to previous loops'' is still too strong;
  \item the original narrower claim remains true at \texttt{trigger\_end}: the final-layer state after the full triggering copy is written is overwhelmingly prompt-focused;
  \item the stricter \texttt{trigger\_start} query is meaningfully different: previous-loop attention is stronger both in the final layer (\texttt{0.054} vs \texttt{0.031}) and in the middle of the stack (\texttt{0.249} vs \texttt{0.193});
  \item but even at \texttt{trigger\_start}, the final layer is still mostly prompt-focused rather than previous-loop-focused (\texttt{prompt mass = 0.576}, \texttt{top1\_prompt = 0.825}).
\end{itemize}

So if the question is literally ``what is the model looking at on the first token of the final repeated copy?'', the answer is: still mostly the prompt in the final layer, with a clearly stronger but still secondary previous-loop signal.

The next still-open refinement is the one-token-earlier \texttt{trigger\_start - 1} query plus matched non-loop controls.

\section*{Future Rollouts}

The saver-side fixes are already on this branch:

\begin{itemize}
  \item prompt-profile archives now save exact \texttt{completion\_token\_ids} plus structured \texttt{loop\_trigger}
  \item the rollout-stats collector now writes \texttt{\_\_rollout\_archive.jsonl.gz} with prompt/completion token IDs and trigger metadata
\end{itemize}

So this exact reconstruction gap should not recur on future rollouts.

\section*{Artifact Bundle}

Raw outputs now live in two sibling directories:

\begin{itemize}
  \item \texttt{outputs/qwen3\_loop\_trigger\_attention\_full\_20260414\_rerun/} for \texttt{trigger\_end}
  \item \texttt{outputs/qwen3\_loop\_trigger\_attention\_trigger\_start\_20260414/} for \texttt{trigger\_start}
\end{itemize}

The updated report PDF stays under the original report bundle:

\begin{itemize}
  \item \texttt{outputs/qwen3\_loop\_trigger\_attention\_full\_20260414\_rerun/qwen3\_loop\_trigger\_attention\_full\_20260414\_rerun.pdf}
\end{itemize}

\end{document}
